\documentclass[a4paper,11pt]{article}
\usepackage[utf8]{inputenc}
\usepackage[T1]{fontenc}
\usepackage[french]{babel}
\usepackage{geometry}
\geometry{top=1.8cm,bottom=1.8cm,left=2.2cm,right=2.2cm}
\usepackage{xcolor}
\usepackage{enumitem}
\usepackage{parskip}
\definecolor{primary}{RGB}{0,82,165}
\definecolor{textdark}{RGB}{18,22,55}
\definecolor{muted}{RGB}{100,100,110}
\definecolor{transition}{RGB}{0,130,80}
\pagestyle{empty}
\setlength{\parindent}{0pt}
\newcommand{\notehead}[2]{%
  \par\vspace{12pt}%
  {\color{white}\colorbox{primary}{\parbox{\dimexpr\linewidth-12pt}{\vspace{3pt}%
    \bfseries\large Diapo #1~--- #2\vspace{3pt}}}}%
  \par\vspace{5pt}%
}
\newcommand{\pivot}[1]{%
  \par\vspace{4pt}%
  {\color{transition}\small\bfseries $\Rightarrow$ #1}%
  \par%
}
\begin{document}
\begin{center}
  {\LARGE\bfseries\color{primary} RansomShield --- Discours de soutenance}\\[4pt]
  {\large\color{textdark} 10 min de pr\'esentation + 5 min de d\'emo}\\[2pt]
  {\small\color{muted} Codjo Fr\'ejus Lo\"ic SANGRONIO --- 24 juin 2026}
\end{center}
\vspace{4pt}
{\color{muted}\small Lire \`a voix haute avant la soutenance. Les \textcolor{transition}{fl\`eches vertes} sont des phrases \`a prononcer, pas \`a lire~: elles enchainent les diapos naturellement. Les \textbf{mots en gras} sont les ancres \`a m\'emoriser en priorit\'e.}
\vspace{2pt}
\hrule

% ─────────────────────────────────────────────────────
\notehead{1/25}{OUVERTURE — 0:00}

Monsieur le Président du jury, Mesdames et Messieurs les membres du jury,
  bonjour.

Je m'appelle Loïc SANGRONIO, et je vais vous présenter aujourd'hui \textbf{RansomShield}~:
  un système de détection et de réponse aux attaques ransomware,
  dans un environnement contrôlé.

Dix minutes de présentation. Cinq minutes de démonstration live,
  sur des machines virtuelles réelles. Et je serai à votre disposition
  pour vos questions à la fin.

\pivot{Commençons par un coup d'œil au plan.}

% ─────────────────────────────────────────────────────
\notehead{2/25}{SOMMAIRE — 0:30}

Cette présentation suit un fil simple~:
  d'abord \textbf{le pourquoi} — le contexte et les objectifs~;
  ensuite \textbf{le quoi} — la conception et les choix techniques~;
  enfin \textbf{le résultat} — les tests, les limites, et la démonstration.

\pivot{Commençons par le contexte.}

% ─────────────────────────────────────────────────────
\notehead{4/25}{CONTEXTE \& PROBLÉMATIQUE — 1:00}

Imaginez~: vos fichiers se renomment les uns après les autres,
  les sauvegardes disparaissent,
  et un message s'affiche à l'écran — «~payez, ou perdez tout~».
  C'est un ransomware.

L'ENISA a mesuré une \textbf{hausse de 73\,\% de ces attaques en 2023}.
  Sans outil dédié, il faut en moyenne \textbf{21 jours pour détecter} l'incident,
  pour un coût qui dépasse souvent le million de dollars.

Le problème fondamental~? Les antivirus classiques détectent par signature.
  Un ransomware inconnu passe à travers.

D'où la question centrale de ce travail~:
  \textbf{comment détecter un comportement ransomware,
  dans un environnement contrôlé, sans déployer de vrai malware~?}

\pivot{C'est précisément ce que RansomShield cherche à faire — voyons les objectifs.}

% ─────────────────────────────────────────────────────
\notehead{5/25}{OBJECTIFS — 2:30}

Six objectifs concrets~:
  surveiller les événements \textbf{en temps réel},
  calculer un \textbf{score de risque},
  générer automatiquement des alertes et des incidents,
  proposer des actions de protection — avec un \textbf{contrôle humain à chaque étape} —,
  sécuriser la console par une \textbf{double authentification},
  et valider l'ensemble sur des \textbf{scénarios d'attaque réalistes}.

\pivot{Pour y répondre, le mémoire s'articule en trois chapitres.}

% ─────────────────────────────────────────────────────
\notehead{6/25}{ORGANISATION DU MÉMOIRE — 3:15}

Revue de littérature. Conception du système. Réalisation et résultats.
  Voyons d'abord les fondements théoriques.

\pivot{Qu'est-ce qu'un ransomware, exactement~?}

% ─────────────────────────────────────────────────────
\notehead{8/25}{DÉFINITION \& COMPORTEMENTS — 3:50}

Le NIST définit le ransomware comme un logiciel malveillant qui
  \textbf{bloque l'accès aux données} — le plus souvent par chiffrement —
  pour ensuite réclamer une rançon.

Ce qui nous intéresse ici~: ces attaques \textbf{ne sont pas silencieuses}.
  Elles laissent des traces avant que les dégâts ne soient irréversibles~:
  des fichiers renommés en \texttt{.locked},
  une note de rançon qui apparaît,
  les sauvegardes qui s'effacent,
  ou des outils système détournés de leur usage normal.

\textbf{Ce sont ces signaux comportementaux qui constituent notre détection.}

\pivot{Mais pourquoi les outils existants ne suffisent-ils pas~?}

% ─────────────────────────────────────────────────────
\notehead{9/25}{SOLUTIONS EXISTANTES — 5:15}

Les antivirus~? Signatures connues~; inefficaces sur les variantes nouvelles.\\
  Les SIEM~? Ils collectent des journaux — sans corrélation comportementale ciblée.\\
  Les EDR commerciaux~? CrowdStrike, SentinelOne sont efficaces — et hors de portée
  pour un contexte académique.

\textbf{RansomShield se positionne dans cet espace vide~:}
  approche comportementale, open-source, avec un opérateur humain
  dans la boucle à chaque décision.

\pivot{Concrètement, comment est-il construit~?}

% ─────────────────────────────────────────────────────
\notehead{11/25}{ARCHITECTURE — 6:15}

Le flux est simple à suivre.
  Les \textbf{agents Python} surveillent les machines
  et envoient les événements à l'\textbf{API Laravel}.
  Le \textbf{moteur de détection} analyse, calcule un score,
  et déclenche une alerte si le seuil est dépassé.
  L'opérateur voit tout sur la \textbf{console SOC}
  et \textbf{valide chaque action avant son exécution}.

Règle d'or du système~: \textbf{aucune action automatique sans validation humaine.}

\pivot{Ce flux repose sur quatre composants — regardons-les.}

% ─────────────────────────────────────────────────────
\notehead{12/25}{LES 4 COMPOSANTS — 7:15}

\textbf{1.~L'agent Python}~: surveille les fichiers en temps réel,
  conserve une file SQLite locale pour ne rien perdre en cas de coupure réseau.\\[3pt]
  \textbf{2.~La console SOC}~: centralise alertes, incidents, et file d'approbation.\\[3pt]
  \textbf{3.~La découverte réseau}~: scanne les sous-réseaux automatiquement,
  mais aucun agent ne s'enrôle sans validation manuelle de l'administrateur.\\[3pt]
  \textbf{4.~Le module de simulation}~: rejoue jusqu'à 22 événements d'attaque
  — sans jamais toucher à un vrai malware.

\pivot{Deux acteurs utilisent ce système — qui sont-ils~?}

% ─────────────────────────────────────────────────────
\notehead{13/25}{DIAGRAMME UML — 8:45}

Deux acteurs, deux rôles distincts.

L'\textbf{administrateur}~: il configure le système,
  gère les agents, et lance les simulations.\\[3pt]
  L'\textbf{analyste SOC}~: il traite l'opérationnel —
  il visualise les alertes, approuve ou rejette les actions,
  consulte la timeline d'audit.
  Mais il n'a pas accès à la configuration.

\pivot{Ces acteurs manipulent des données — voyons comment elles sont structurées.}

% ─────────────────────────────────────────────────────
\notehead{14/25}{MODÈLE DE DONNÉES — 9:45}

Huit tables MySQL. Le pipeline se lit de gauche à droite~:
  \textbf{agent $\to$ événement $\to$ alerte $\to$ incident $\to$ action $\to$ journal d'audit.}

Chaque action est signée~: on sait \textbf{qui} l'a validée, et \textbf{quand}.
  Aucune décision n'est anonyme.

\pivot{Passons maintenant à la réalisation concrète.}

% ─────────────────────────────────────────────────────
\notehead{16/25}{CHOIX TECHNIQUES — 11:00}

Backend~: \textbf{Laravel 11} et PHP 8.3, base \textbf{MySQL},
  API sécurisée par une clé propre à chaque agent.\\[3pt]
  Agent~: \textbf{Python + Watchdog}, file locale \textbf{SQLite}.\\[3pt]
  Frontend~: \textbf{Chart.js} en local, sans dépendance externe.\\[3pt]
  Infrastructure de test~: \textbf{3 VMs KVM}, double authentification activée.

\pivot{Concrètement, comment l'agent fonctionne-t-il au démarrage~?}

% ─────────────────────────────────────────────────────
\notehead{17/25}{AGENT \& MOTEUR — 12:15}

Au premier démarrage, l'agent s'enrôle via un \textbf{token à usage unique}
  — il obtient une clé API permanente.
  Ensuite~: surveillance des dossiers avec Watchdog,
  et \textbf{heartbeat toutes les 30 secondes}.

Le moteur calcule un score cumulatif.
  Exemple réel~: fichier renommé \texttt{.locked} = \textbf{80 points},
  note de rançon détectée = \textbf{55 points},
  total \textbf{135} $\to$ \textbf{alerte critique + incident créé automatiquement}.
  Durée totale~: \textbf{moins de 5 secondes.}

\pivot{Côté analyste, qu'est-ce que la console offre~?}

% ─────────────────────────────────────────────────────
\notehead{18/25}{CONSOLE SOC — 13:30}

Dix pages fonctionnelles.
  Le dashboard se rafraîchit \textbf{toutes les 30 secondes} automatiquement.

Chaque action de protection passe par une \textbf{file d'approbation}~:
  l'analyste valide ou rejette avant toute exécution.
  Les notifications partent sur \textbf{quatre canaux simultanés}~:
  web, e-mail, son, et webhook.

\textbf{Vous allez voir tout ça dans quelques instants.}

\pivot{Mais d'abord — qu'est-ce que les tests ont donné~?}

% ─────────────────────────────────────────────────────
\notehead{19/25}{TESTS \& RÉSULTATS — 14:45}

Cinq scénarios d'attaque, de 7 à 22 événements.
  \textbf{Détection en moins de 5 secondes} à chaque fois.
  Zéro perte d'événement.
  Pipeline validé \textbf{de bout en bout}~:
  de l'événement jusqu'à la timeline d'audit.

Une limite que j'assume honnêtement~: les tests ont été conduits
  \textbf{uniquement sur Linux}.
  Un attaquant qui ralentirait délibérément son rythme resterait moins visible.

\pivot{Maintenant, place à la démonstration.}

% ─────────────────────────────────────────────────────
\notehead{20/25}{DÉMO LIVE — 15:00 \hspace{6pt}{\normalsize\normalfont(checklist — 5 min, hors chrono)}}

\begin{enumerate}[leftmargin=1.4em,itemsep=3pt]
\item Ouvrir \texttt{http://10.20.0.1} — console SOC
\item Se connecter~: identifiants + \textbf{code 2FA}
\item Vérifier~: \textbf{3 agents en ligne} (heartbeat vert)
\item \textbf{Agents} $\to$ VM-02 $\to$ \textbf{Lancer simulation} $\to$ «~Kill chain complète~» (22 évts)
\item \textbf{Dashboard}~: le score monte en temps réel
\item \textbf{Alertes}~: niveau Critique + signaux détectés
\item \textbf{Incidents}~: ouvrir l'incident auto-créé
\item \textbf{File d'approbation}~: approuver l'isolation réseau
\item \textbf{Timeline}~: tout est horodaté et tracé
\end{enumerate}
{\color{muted}\textit{Rester calme~; chaque étape a le temps qu'il lui faut.}}

\pivot{Revenons à la présentation pour conclure.}

% ─────────────────────────────────────────────────────
\newpage
\notehead{22/25}{CONCLUSION \& PERSPECTIVES — 20:00}

RansomShield couvre \textbf{tout le cycle de traitement d'un incident}~:
  collecte, analyse, alerte, réponse, audit.
  Agent multi-plateforme. Console complète.
  Moteur réactif. \textbf{Contrôle humain systématique.}

Ce travail a ses limites~: environnement contrôlé,
  réponses partiellement simulées,
  tests conduits sur Linux uniquement.

Les perspectives sont claires~:
  d'abord \textbf{le chiffrement HTTPS},
  puis des rôles utilisateurs distincts,
  puis à terme du \textbf{machine learning} pour affiner la détection,
  et une \textbf{intégration avec les SIEM} du marché.

Je vous remercie pour votre attention.
  \textbf{Je suis à votre disposition pour vos questions.}

\vspace{16pt}
% ─────────────────────────────────────────────────────
\notehead{23/25}{QUESTIONS DU JURY — réponses prêtes}

\textbf{Pourquoi Python et pas C ou Go~?}\\
  \hspace*{1.2em}Watchdog exploite \texttt{inotify} nativement — CPU minimal,
  et portable Linux, Windows, macOS sans recompilation.\\[6pt]

\textbf{Vous gérez les faux positifs comment~?}\\
  \hspace*{1.2em}Les seuils sont configurables. Et surtout~:
  \textbf{l'opérateur humain valide ou rejette chaque action avant qu'elle ne s'exécute} —
  c'est la garantie centrale du système.\\[6pt]

\textbf{Pourquoi pas HTTPS~?}\\
  \hspace*{1.2em}Réseau LAN contrôlé pour ce projet.
  Le chiffrement HTTPS est la première perspective pour un passage en production.\\[6pt]

\textbf{Quelle différence avec un antivirus~?}\\
  \hspace*{1.2em}Un antivirus cherche des signatures connues.
  RansomShield analyse le \textbf{comportement}~:
  il détecte même une variante jamais vue auparavant.\\[6pt]

\textbf{Testé sur Windows~?}\\
  \hspace*{1.2em}Non — limite assumée dans ce travail. Linux uniquement.
  C'est une perspective d'extension explicite.\\[6pt]

\textbf{La solution passe-t-elle à l'échelle~?}\\
  \hspace*{1.2em}Testé sur 3 VMs. L'API Laravel peut monter en charge,
  mais des tests de performance à grande échelle restent à conduire.\\[8pt]

{\color{muted}\textit{Reformuler la question avant de répondre. Garder les captures d'écran prêtes.}}

\end{document}
